\section{Stage I: what can identify or contain the support?}
\label{sec:two-stage-discovery}

The two-stage theorem deliberately parameterizes discovery by
$W_{\rm disc}$.  The following table records which engines currently close
that parameter and which merely expose the remaining interface.

\begin{center}
\small
\begin{tabularx}{\textwidth}{@{}p{0.20\textwidth}p{0.20\textwidth}p{0.22\textwidth}X@{}}
\toprule
Stage~I engine & Returned object & Charged work & Status and limitation \\
\midrule
Classical APPR & numerical point and envelope &
$\Theta(1/(\alpha\rho))$ worst case &
unconditional direct return; tail is optional \\
Hard-capped Green ball & set-only approximate envelope &
$O(B)$ screen, then fixed tail &
unconditional adaptive lane; may abort to APPR \\
PPR leakage screen & coarse principal point plus boundary error bars &
proposal work plus $\widetilde O(\vol(E)/\sqrt\alpha)$ &
direct local gate; no RPPR support or zero-key decisions \\
Exact tree response & exact RPPR point and face &
$\widetilde O((1+\vol(S^\star))/\sqrt\alpha)$ &
unconditional on promised point-source trees \\
Unicyclic/cactus response & exact face &
product-scale structural bounds &
exact-real; scope depends on the reporter theorem \\
Hard-capped dynamic obstacle & point/face or retained envelope &
$W_{\rm DS}(2/\rho)+O(1/\rho)$ &
conditional on an exposure-local response/reporter oracle \\
Restarted AESP-CD obstacle batches & point and envelope from a gap &
$\widetilde O(\vol(U)/\sqrt\alpha)$ per fixed envelope &
the solve is closed; serial discovery/reporting is open \\
Fixed-envelope AESP mass screen & lower-safe PPR checkpoint &
$\widetilde O(\vol(E)/\sqrt\alpha)$ per tested envelope &
no support needed; hard-capped and instance-selective by the star barrier \\
Monotone GS/SOR with $\omega=1$ & lower-safe numerical state &
instance/gap dependent &
robust verifier or fallback; no universal accelerated discovery rate \\
\bottomrule
\end{tabularx}
\end{center}

\subsection{When Stage I can certify the exact support}

Exact support remains a valid optional handoff when the instance has a
visible KKT margin.  The following interval test makes the quantifier precise
and avoids solving a fixed face to symbolic exactness.

\begin{proposition}[Residual-error exact-face verifier]
\label{prop:two-stage-interval-face-verifier}
Let $A\ni v$ be connected, let
$p_A=Q_A^{-1}c_{\rho,A}$, and let $\widehat p_A$ satisfy
\begin{equation}
 \max_{i\in A}
 \frac{|(p_A)_i-(\widehat p_A)_i|}{\sqrt{d_i}}
 \leq\eta.
 \label{eq:two-stage-face-error-bar}
\end{equation}
If
\begin{align}
 (\widehat p_A)_i-\eta\sqrt{d_i}&>0 &&(i\in A),
 \label{eq:two-stage-face-primal-interval}\\
 c_{\rho,j}-Q_{jA}\widehat p_A
 +\frac{1-\alpha}{2}
  \frac{\deg_A(j)}{\sqrt{d_j}}\eta&\leq0
 &&(j\in\partial A),
 \label{eq:two-stage-face-dual-interval}
\end{align}
then $A=S^\star(\rho)$.  The residual choice
\begin{equation*}
 \eta=\frac{\norm{c_{\rho,A}-Q_A\widehat p_A}_2}{\alpha}
\end{equation*}
is sufficient.  Under positive primal and strict negative boundary margins,
fixed-envelope Chebyshev/AESP reaches this certificate in
$\widetilde O(\vol(A)/\sqrt\alpha)$ work, with the inverse margin entering
only logarithmically.
\end{proposition}

\begin{proof}
The first interval gives $p_A>0$.  For a boundary row, the error contribution
to its exact key is at most
\begin{equation*}
 \sum_{i\in A}(-Q_{ji})\eta\sqrt{d_i}
 =\frac{1-\alpha}{2}
   \frac{\deg_A(j)}{\sqrt{d_j}}\eta.
\end{equation*}
Thus the second interval proves
$c_{\rho,j}-Q_{jA}p_A\leq0$.  A nonboundary exterior row has no active
neighbor and has strictly negative nonseed load.  The zero extension of
$p_A$ therefore satisfies the global obstacle KKT conditions, proving
$A=S^\star$.  Strong convexity supplies the residual error bar.  Geometric
residual reduction on the fixed spectrum $[\alpha,1]$ gives the work claim
when the declared inequalities have strict margins.
\end{proof}

This is the clean two-stage ``identify $S^\star$, then restart'' theorem, but
its scope should not be overstated.  It verifies a proposed face cheaply; it
does not expose that face on an arbitrary graph.  With a zero active or dual
margin, no uniform finite-precision procedure can certify the exact support
in finite work.  The boundary-leakage and approximate-envelope exits are
therefore the margin-free main contracts.

\begin{theorem}[Interruptible APPR-to-exact-face race]
\label{thm:two-stage-appr-face-race}
Choose $0<\rho\leq\epsppr$ and a face-solve error budget
$\eta_{\max}\leq\epsppr-\rho$, and run priority APPR toward its ordinary
numerical stopping rule at target $\epsppr$.  At any charged checkpoint, a proposal
lane may snapshot the current APPR support or a charged local expansion $A$
and apply \cref{prop:two-stage-interval-face-verifier}.  If the verifier
passes with error bar $\eta\leq\eta_{\max}$, return its zero-padded point.  Race
the cumulative verifier work one-for-one against the continuing APPR state;
if APPR finishes first, abort verification and return APPR.  Then the output
always satisfies the semantic $\epsppr$ target and the total work is at most
twice direct priority-APPR work.  A successful exact structural solve has
$\eta=0$, permits the tuned choice $\rho=\epsppr$, identifies
$S^\star(\rho)$, and constitutes a certified screen-or-return exit.  A
strict two-stage variant may instead retain only the face certificate and
solve $Q_Au_A=b_A$; if that additional work is included in the proposal lane,
the same fair-race guarantee applies.
\end{theorem}

\begin{proof}
On a passing face, the interval KKT test proves that the exact restricted
point is $x_\rho^\star$.  The RPPR--PPR comparison contributes at most
$\rho$ normalized error and the certified face error contributes at most
$\eta\leq\eta_{\max}$, whose sum is at most $\epsppr$.  The APPR exit has its standard
$\epsppr$ certificate.  Until either exit, equal scheduling gives the
verifier no more charged work than the direct lane, proving the factor-two
bound.  Snapshotting does not mutate the continuing APPR state.
\end{proof}

The deterministic implementation checkpoints at residual parameters
$16\rho,8\rho,4\rho,2\rho$, tries the APPR support and its one-hop expansion,
uses $\rho=\epsppr$, and applies the implemented leaf/unicyclic solver for the exact-face test.  Of
$90$ path, cycle, star, binary-tree, broom, and grid cases on
$1024$--$2047$ vertices, $19$ nontrivial verifier traces finish before the
direct lane and return exact support; four exits strictly beat direct APPR.
Their ratios are $0.132$ and $0.276$ on stars and $0.590$ and $0.921$ on
binary trees.  Those counts return the exact RPPR verifier point immediately.
Charging a separate fresh ordinary-PPR solve after certification leaves four
standalone wins and two fair-race wins, both on stars.  Every failed
trace is charged at most twice direct work.  These numbers show a genuine
but margin- and structure-selective exact-support lane, not a universal
discovery theorem.
Retaining the structural factors and charging only the second backsolve raises
the standalone count to six while leaving the fair count at two.

\begin{corollary}[Positive-residual SOR discovers, one face solve finishes]
\label{cor:two-stage-sor-support-then-solve}
Start from the zero lower checkpoint for the RPPR problem at
$\rho=\epsppr$ when the fixed-face solve is exact.  Priority SOR may add every currently positive residual row
to a connected candidate $A$, perform only lower-safe coordinate updates,
and launch the exact-face verifier at geometrically increasing charged-work
checkpoints.  Race this lane against direct APPR.  Every proposed coordinate
lies in $S^\star(\rho)$; a passing verifier therefore occurs exactly at the
full support.  Its RPPR point may return immediately, or the certificate may
trigger one ordinary-PPR solve with the same matrix.  The combined
algorithm is always correct and costs at most twice the faster completed
lane.  No claim of a graph-universal accelerated discovery time is implied.
\end{corollary}

\begin{proof}
Let $z\leq x_\rho^\star$ be the current lower-safe point and let $i$ be an
inactive row.  If $(c_\rho-Qz)_i>0$ but $(x_\rho^\star)_i=0$, then the
Stieltjes signs and $z\leq x_\rho^\star$ give
$(c_\rho-Qz)_i\leq(c_\rho-Qx_\rho^\star)_i\leq0$, a contradiction.  Thus
positive residual rows are support-safe.  The interval-face proposition and
the same fair-race argument as in
\cref{thm:two-stage-appr-face-race} prove the remaining statements.
\end{proof}

The implementation checks at powers of two in accumulated SOR degree work,
uses no high-accuracy changing-face solve, and performs the implemented
tree/unicyclic solve only on snapshots.  In the same $90$ cases, this exact
lane finishes before direct APPR in eight cases.  After charging both
fair-race lanes, five remain strict end-to-end wins: four stars and the
low-$\alpha$ binary tree, with ratios $0.088$, $0.118$, $0.345$, $0.444$,
and $0.593$.  This is the
closest executable realization of ``AESP-CD/SOR discovers $S^\star$, then
restart once'' in the note: even after a separate fresh ordinary-PPR solve is
charged, the same five cases remain fair-race wins.  Paths and cycles expose the known serial-frontier
limitation and correctly fall back to APPR.
The retained-factor charge raises the standalone count from five to seven and
leaves the fair count at five.

The charged diagnostic reinforces this distinction.  On five graphs with
$1024$--$2047$ vertices, $\epsppr=10^{-4}$, and
$\rho=\epsppr/2$, literal monotone RPPR Gauss--Seidel often returns a smaller
set than the Green envelope, but its median discovery work is $7.92$ times
Green--CG and $9.26$ times direct PPR.  At $\alpha=0.01$, the 1024-path exact
support has only $41$ vertices and volume $81$, yet discovering it costs
$45518$ versus $5565$ for Green--CG and $6312$ for direct PPR.  On the
1024-star the corresponding works are $1054713$, $6138$, and $116622$.
These figures are not a lower bound against every accelerated active-set
method.  They do show why a smaller exact face is not by itself a useful
Stage~I contract: changing-face work must be paid before its smaller Stage~II
can help.

\subsection{Potential-to-envelope conversion: the direct AESP/SOR handoff}

The approximate-envelope interface creates a handoff that does not need a
final boundary KKT scan.

\begin{proposition}[A certified RPPR gap yields an output-sized envelope]
\label{prop:two-stage-gap-to-envelope}
Let $z$ be a feasible RPPR point with a certified gap
\begin{equation}
 F_\rho(z)-F_\rho(x^\star(\rho))
 \leq\frac{\alpha\delta^2}{8}.
 \label{eq:two-stage-gap-certificate}
\end{equation}
Define
\begin{equation}
 E_\delta(z):=
 \left\{i:\frac{|z_i|}{\sqrt{d_i}}>\frac\delta2\right\}.
 \label{eq:two-stage-threshold-envelope}
\end{equation}
Then $E_\delta(z)$ is a $\delta$-approximate-envelope certificate.  If
$\norm{D^{1/2}z}_1\leq M$, then
\begin{equation}
 \vol(E_\delta(z))<\frac{2M}{\delta}.
 \label{eq:two-stage-threshold-envelope-volume}
\end{equation}
In particular, a lower-safe point $0\leq z\leq x^\star(\rho)$ has $M\leq1$.
\end{proposition}

\begin{proof}
Strong convexity and \cref{eq:two-stage-gap-certificate} give
$\norm{z-x^\star}_2\leq\delta/2$.  If $i\notin E_\delta(z)$, then
\begin{equation*}
 \frac{x_i^\star}{\sqrt{d_i}}
 \leq\frac{|x_i^\star-z_i|}{\sqrt{d_i}}
     +\frac{|z_i|}{\sqrt{d_i}}
 \leq\delta,
\end{equation*}
because $d_i\geq1$.  Every $i\in E_\delta(z)$ contributes more than
$\delta d_i/2$ to $\norm{D^{1/2}z}_1$, proving the volume bound.  For a
lower-safe point, Stieltjes monotonicity gives
$z\leq x^\star\leq x^0$, while the point-source PPR mass identity is
$\langle D^{1/2}\one,x^0\rangle=1$.
\end{proof}

Taking $\rho=\delta=\epsppr/2$ shows exactly what AESP-CD or monotone
SOR must accomplish for an early return: produce a sparse feasible state and
certify gap at most $\alpha\epsppr^2/32$.  Then
$\vol(E_\delta)<4/\epsppr$, and
\cref{cor:two-stage-numerical-handoff-returns} returns the thresholded point
without a linear tail.  If an application insists on a canonical restricted
linear output, the three-way split and
\cref{cor:two-stage-three-budget} remain available.
This is strictly weaker than identifying $S^\star$ or certifying every zero
key.  The unresolved issue is the charged work needed to obtain the gap on an
arbitrary locally exposed graph.

\paragraph{Mass-capped accelerated AESP lane.}
Coordinatewise lower safety is sufficient for the mass bound, but not
necessary.  The RPPR optimum belongs to
\begin{equation}
 \mathcal C_1:=
 \{z\geq0:\langle D^{1/2}\one,z\rangle\leq1\}.
 \label{eq:two-stage-mass-cap}
\end{equation}
Therefore accelerated projected gradient may use scratch iterates projected
onto $\mathcal C_1$ without changing the optimum.  Once its certified gap
reaches \cref{eq:two-stage-gap-certificate}, thresholding gives volume below
$2/\delta$ even if the scratch was not coordinatewise lower-safe.  The outer
iteration count is
\begin{equation*}
 O\!\left(\alpha^{-1/2}
   \log\frac{C}{\alpha\delta^2}\right).
\end{equation*}
Thus, if every charged scratch matrix application stays in an exposed
envelope of volume $O(1/\delta)$, this lane has the target
$\widetilde O(1/(\delta\sqrt\alpha))$ Stage-I work.  This is a conditional
locality theorem, not a universal one: a small mass on a newly reached
high-degree coordinate can still make a matrix application expensive.  The
mass cap removes the support/publication obstruction and leaves precisely the
trajectory-exposure obstruction.

The cap does give a graph-uniform one-sided exposure fact.  Write
$a=(1+\alpha)/2$, $b=(1-\alpha)/2$, and
$N=D^{-1/2}AD^{-1/2}$.

\begin{proposition}[Mass cap controls every positive-violation batch]
\label{prop:two-stage-mass-cap-positive-exposure}
For $z\in\mathcal C_1$, put $r=c_\rho-Qz$ and
\begin{equation*}
 C_+(z):=\norm{D^{1/2}[r]_+}_1,
 \qquad
 P_\eta(z):=\left\{i:\frac{[r_i]_+}{\sqrt{d_i}}\geq\eta\right\}.
\end{equation*}
Then
\begin{equation}
 C_+(z)\leq\alpha+b\leq1,
 \qquad
 \vol(P_\eta(z))\leq\frac{\alpha+b}{\eta}.
 \label{eq:two-stage-positive-batch-volume}
\end{equation}
Moreover, a nonseed coordinate with $r_i>0$ has
$d_i<b/(\alpha\rho)$; if the seed residual is positive, then
$d_v<(\alpha+b)/(\alpha\rho)$.
\end{proposition}

\begin{proof}
The source and $Nz$ are nonnegative, whereas the regularizer and diagonal
parts of $r$ are nonpositive.  Hence
\begin{equation*}
 [r]_+\leq \alpha D^{-1/2}e_v+bNz.
\end{equation*}
The normalized adjacency preserves degree-weighted mass:
$\langle D^{1/2}\one,Nz\rangle=\langle D^{1/2}\one,z\rangle$.
This proves the first inequality.  Every coordinate in $P_\eta$ contributes
at least $\eta d_i$ to $C_+$, proving the volume bound.  Finally, for a
nonseed coordinate, multiplication by $\sqrt{d_i}$ and omission of the
nonpositive diagonal term give
\begin{equation*}
 \sqrt{d_i}r_i
 \leq-\alpha\rho d_i
     +b\sum_{j\sim i}\frac{z_j}{\sqrt{d_j}}
 \leq-\alpha\rho d_i+b.
\end{equation*}
Positivity yields the degree screen.  At the seed, adding the weighted source
contribution $\alpha$ gives the stated weaker bound.
\end{proof}

For a shifted AESP subproblem centered at another point $y\in\mathcal C_1$,
the same proof bounds the positive residual mass by $\alpha+b+\kappa$ and a
positive-violation batch by $(\alpha+b+\kappa)/\eta$.  This is useful but not
yet the missing theorem: signed scratch can create negative KKT keys on its
positive support, and repeated outer restarts are not paid by
\cref{eq:two-stage-positive-batch-volume}.  The proposition proves that the
obstruction is not a single enormous positive batch; it is signed trajectory
maintenance and repeated exposure.

For a lower-safe publication, even the objective gap is unnecessary.  Define
the maximum positive key density
\begin{equation}
 \eta_+(z):=\norm{D^{-1/2}[c_\rho-Qz]_+}_\infty.
 \label{eq:two-stage-positive-key-density}
\end{equation}

\begin{proposition}[One-heap lower-safe envelope handoff]
\label{prop:two-stage-residual-to-envelope}
If $0\leq z\leq x^\star(\rho)$ and
\begin{equation}
 \eta_+(z)\leq\frac{\alpha\delta}{2},
 \label{eq:two-stage-residual-handoff}
\end{equation}
then
\begin{equation}
 0\leq x^\star(\rho)-z
 \leq\frac\delta2D^{1/2}\one.
 \label{eq:two-stage-lower-safe-error-box}
\end{equation}
Consequently $E_\delta(z)$ is a $\delta$-approximate envelope of volume below
$2/\delta$.
\end{proposition}

\begin{proof}
Put $e=x^\star-z\geq0$ and $r=c_\rho-Qz$.  On an active optimum coordinate,
complementarity gives $(Qe)_i=r_i$.  On an inactive optimum coordinate,
$e_i=z_i=0$ and the Stieltjes signs give $(Qe)_i\leq0$.  Therefore
\begin{equation*}
 Qe\leq[r]_+
 \leq\frac{\alpha\delta}{2}D^{1/2}\one
 =Q\left(\frac\delta2D^{1/2}\one\right).
\end{equation*}
Inverse positivity proves \cref{eq:two-stage-lower-safe-error-box}.  Outside
$E_\delta(z)$, both $z_i/\sqrt{d_i}\leq\delta/2$ and
$e_i/\sqrt{d_i}\leq\delta/2$, proving the envelope claim.  Lower safety gives
$\norm{D^{1/2}z}_1\leq1$, so
\cref{eq:two-stage-threshold-envelope-volume} gives the volume bound.
\end{proof}

\begin{corollary}[A numerical discovery certificate can return early]
\label{cor:two-stage-numerical-handoff-returns}
Suppose either the gap certificate of
\cref{prop:two-stage-gap-to-envelope} or the lower-safe key certificate of
\cref{prop:two-stage-residual-to-envelope} holds.  Zero every coordinate of
$z$ outside $E_\delta(z)$ and take the positive part; call the resulting
vector $\bar z$.  Then
\begin{equation}
 \norm{D^{-1/2}(\bar z-x^0)}_\infty\leq\rho+\delta.
 \label{eq:two-stage-numerical-early-output}
\end{equation}
Thus if $\rho+\delta\leq\epsppr$, this lane should return $\bar z$ without
running Stage~II.  Likewise an exact RPPR face solution is already
$\rho$-accurate for PPR.  The fixed linear tail is needed only for a set-only
certificate, for a prescribed canonical linear output, or as optional
refinement.
\end{corollary}

\begin{proof}
In the gap case, strong convexity gives
$\norm{z-x^\star}_2\leq\delta/2$; taking the positive part cannot increase
distance to $x^\star\geq0$.  Inside $E_\delta$ the normalized error is at most
$\delta/2$, and outside it the proof of
\cref{prop:two-stage-gap-to-envelope} gives
$x_i^\star/\sqrt{d_i}\leq\delta$.  The one-heap case has the same inside and
outside bounds by \cref{eq:two-stage-lower-safe-error-box}.  In either case
$\bar z$ is $\delta$-accurate for RPPR; combine with
\cref{eq:two-stage-rppr-bridge}.
\end{proof}

For a zero-start monotone obstacle Gauss--Seidel trace, lower safety is an
invariant.  A lazy heap of positive normalized keys exposes both the next
priority update and the scalar certificate \cref{eq:two-stage-residual-handoff}.
The remaining question is not observability of the handoff; it is whether the
work needed to reach the $\alpha\delta$ key scale can be bounded by the
accelerated target.

There is nevertheless a clean unconditional nonaccelerated bound.  Put
$a=(1+\alpha)/2$, $b=(1-\alpha)/2$, $r=c_\rho-Qz$, and define the positive
residual mass
\begin{equation*}
 C_+(z):=\sum_i\sqrt{d_i}[r_i]_+.
\end{equation*}

\begin{theorem}[Priority monotone-SOR Stage~I]
\label{thm:two-stage-priority-sor-discovery}
Start from $z=0$.  While
$\eta_+(z)>\alpha\delta/2$, choose a positive-residual coordinate
maximizing $r_i^2/d_i$ and make the exact Gauss--Seidel update
\begin{equation}
 z_i^+=z_i+\frac{r_i}{a}.
 \label{eq:two-stage-priority-sor-update}
\end{equation}
Maintain changed neighbor residuals in a lazy maximum heap.  The method is
lower-safe, terminates with the certificate of
\cref{prop:two-stage-residual-to-envelope}, and uses degree work
\begin{equation}
 \widetilde O\!\left[
  \frac1\alpha\min\left\{
   \frac{2}{\delta},
   \frac1\rho\left(1+\log_+\frac{2}{\delta}\right)
  \right\}
 \right].
 \label{eq:two-stage-priority-sor-work}
\end{equation}
With $\rho=\delta=\epsppr/2$, this gives a complete unconditional
point-source PPR solver with $\widetilde O(1/(\alpha\epsppr))$ work.
By \cref{cor:two-stage-numerical-handoff-returns}, it should stop after
Stage~I; appending the accelerated linear solve is correct but unnecessary.
\end{theorem}

\begin{proof}
Inductively $0\leq z\leq x^\star$.  Indeed, if $r_i>0$, then $i\in S^\star$;
otherwise $x_i^\star=z_i=0$ and the Stieltjes signs would give $r_i\leq0$.
Using the active optimality equation and $x^\star-z\geq0$ shows that the
coordinate step in \cref{eq:two-stage-priority-sor-update} is at most
$x_i^\star-z_i$.  The update sets $r_i$ to zero and only increases neighbor
residuals.

Let $B_+=\sum_{i:r_i>0}d_i$.  Positive residuals are support-safe, whence
$B_+\leq\vol(S^\star)\leq1/\rho$.  The priority coordinate obeys
\begin{equation*}
 \frac{r_i}{\sqrt{d_i}}\geq\frac{C_+(z)}{B_+}.
\end{equation*}
The update sets $r_i$ to zero and increases a neighbor $j$ by
$br_i/(a\sqrt{d_id_j})$.  Since the positive-part map is one-Lipschitz and
$i$ has $d_i$ neighbors,
\begin{equation*}
 C_+(z^+)
 \leq C_+(z)-\sqrt{d_i}r_i
       +\frac{b}{a}\sqrt{d_i}r_i
 \leq C_+(z)\left(1-\frac{\alpha\rho}{a}d_i\right).
\end{equation*}
Here $1-b/a=\alpha/a$; the factor is nonnegative because
$d_i\leq B_+\leq1/\rho$ and $\alpha\leq a$.  Multiplication over the degree work gives the direct
residual-mass contraction
\begin{equation*}
 C_+(z)\leq C_+(0)
 \exp\!\left(-\frac{\alpha\rho}{a}W\right)
 \leq\alpha\exp\!\left(-\frac{\alpha\rho}{a}W\right).
\end{equation*}
The last inequality follows by evaluating the sole possible positive residual
at the point source.  Since $\eta_+(z)\leq C_+(z)$, the stopping test is
certainly true after degree work at most
\begin{equation*}
 O\!\left[
  \frac{a}{\alpha\rho}
  \left(1+\log_+\frac{2}{\delta}\right)
 \right].
\end{equation*}
There is also a direct additive bound.  Before termination the chosen key
density exceeds $\alpha\delta/2$, so the first inequality in the mass update
shows
\begin{equation*}
 C_+(z)-C_+(z^+)>
 \frac{\alpha^2\delta}{2a}d_i.
\end{equation*}
Since $C_+(0)\leq\alpha$, summing until termination gives
$W<2a/(\alpha\delta)$.  Taking the smaller of the exponential and additive
bounds proves \cref{eq:two-stage-priority-sor-work}.
Heap changes add only logarithmic factors.  At termination
\cref{prop:two-stage-residual-to-envelope} supplies the envelope and
\cref{cor:two-stage-numerical-handoff-returns} supplies the early output.
\end{proof}

The same heap also gives a direct, unregularized PPR baseline.

\begin{corollary}[Direct priority-SOR PPR]
\label{cor:two-stage-direct-priority-ppr}
Set $\rho=0$ and run the priority rule until
$\eta_+(z)\leq\alpha\epsppr$.  Then
\begin{equation*}
 0\leq x^0-z\leq\epsppr D^{1/2}\one,
\end{equation*}
so $z$ itself meets the semantic PPR target.  The degree work is below
$a/(\alpha\epsppr)$ up to heap logarithms.
\end{corollary}

\begin{proof}
Apply \cref{prop:two-stage-residual-to-envelope} with $\rho=0$ and
$\delta=2\epsppr$, but do not threshold the numerical vector.  The error-box
conclusion is exactly the displayed bound.  The additive part of the proof of
\cref{thm:two-stage-priority-sor-discovery} gives
$W<2a/(\alpha\delta)=a/(\alpha\epsppr)$.
\end{proof}

This theorem matches the classical $1/\alpha$ discovery scale rather than
the desired $1/\sqrt\alpha$ scale.  Its value is architectural: it proves a
complete AESP/SOR-compatible certificate baseline without APPR support
sandwiches or exact zero decisions.  It also exposes a decisive redundancy:
at this work scale the numerical discovery vector has already solved the PPR
target, so Stage~II cannot improve the worst-case bound.  Improving
\cref{eq:two-stage-priority-sor-work} requires accelerated safe progress or a
strictly cheaper set-only screening certificate, not a different terminal
solver.

\subsection{A genuine set-only lane: the hard-capped Green ball}

The point-source assumption permits a screen that computes no obstacle
values at all.  The case $\alpha=1$ is the immediate one-coordinate solve;
for $0<\alpha<1$, put
\begin{equation*}
 \lambda_\alpha:=\frac{1-\sqrt\alpha}{1+\sqrt\alpha},
 \qquad
 R_\delta:=
 \left\lceil
  \frac{\log(2/\delta)}
       {\log((1+\sqrt\alpha)/(1-\sqrt\alpha))}
 \right\rceil,
\end{equation*}
and let $B(v,R_\delta-1)$ be the graph ball around the seed (with radius zero
if the displayed integer is zero).

\begin{theorem}[Hard-capped Green-ball screen]
\label{thm:two-stage-green-ball-screen}
For $0<\alpha<1$, the set
$E=B(v,R_\delta-1)$ is a $\delta$-approximate envelope for the unregularized
PPR solution ($\rho=0$).  Consequently a certified adjacency-list BFS with
hard degree-volume cap $B$ has the following behavior:
\begin{enumerate}[label=(\roman*),leftmargin=2.2em]
\item if the whole ball is materialized before the cap, a fixed principal
  solve to tolerance $\tau$ returns $(\delta+\tau)$-accurate PPR using
  \begin{equation*}
   O\!\left(B+\frac{B}{\sqrt\alpha}
       \log\frac{C}{\tau}\right)
  \end{equation*}
  charged work;
\item if the next retained row would exceed the cap, the lane aborts after
  $O(B)$ work without weakening a fallback solver.
\end{enumerate}
With $B=2/\epsppr$ and $\delta=\tau=\epsppr/2$, a successful lane has the
aspirational product work.  On failure, switching to direct APPR gives the
unconditional bound $O(1/(\alpha\epsppr))$ plus the lower-order failed-screen
budget.  More sharply, alternate charged work quanta of this lane and direct
APPR.  If the ball and fixed solve finish, the combined work is at most twice
the faster lane; if the cap aborts, remove the screen and continue APPR, for
total work at most twice direct APPR.  Thus the screen is an unconditional
adaptive algorithm even when it is not a universal product-work theorem.
\end{theorem}

\begin{proof}
The degree-$k$ Chebyshev approximation to $Q^{-1}$ and graph sparsity give the
point-source Green bound
\begin{equation*}
 0\leq x_i^0\leq2\lambda_\alpha^{\operatorname{dist}(v,i)}.
\end{equation*}
This is the provider lemma \path{lem:graph-support-radius}, before its RPPR
threshold argument.  If $i\notin E$, then
$\operatorname{dist}(v,i)\geq R_\delta$, hence
$x_i^0/\sqrt{d_i}\leq2\lambda_\alpha^{R_\delta}\leq\delta$.
Apply \cref{prop:two-stage-envelope-linearization} with $\rho=0$ (or its
identical Schur-complement proof) and then the approximate-envelope
composition with terminal tolerance $\tau$.

During BFS, query a candidate degree before scanning its row.  Every retained
row and every discovered candidate is charged to an already scanned
incidence, so a failed cap costs $O(B)$; a successful materialization costs
$O(\vol(E))\leq O(B)$.  The two displayed branch bounds follow.  Since
$B=2/\epsppr\leq2/(\alpha\epsppr)$, the failed attempt changes the APPR
fallback only by a constant factor.  The fair-race statement follows by
giving each live lane one work quantum in turn: before either lane receives
$W$ quanta, fewer than $2W+O(1)$ total quanta have elapsed.  A failed screen
has consumed no more quanta than the concurrently advanced APPR lane.
\end{proof}

This is the first concrete lane in the note for which Stage~II is not
logically redundant: Stage~I returns only graph coordinates.  It is not a
general-graph accelerated theorem, because a high-growth ball may hit the cap
immediately.  That failure is observable and safe.  The deterministic sweep
finds successful low-growth cases, but direct APPR is still faster on every
tested coarse equal-accuracy case.  The high-accuracy sweep below changes the
empirical conclusion while preserving the same theorem and implementation.

\begin{corollary}[Structural Green tail]
\label{cor:two-stage-green-structural-tail}
In the successful branch of
\cref{thm:two-stage-green-ball-screen}, suppose the exposed principal graph
has a certified elimination order of width at most $w$.  Replacing the
iterative tail by sparse $LDL^\top$ elimination and back substitution returns
the exact principal PPR point in
\begin{equation}
 O\bigl(B+w^3B\bigr)
 \label{eq:two-stage-green-structural-work}
\end{equation}
exact-real work and $O(w^2B)$ state.  In particular, on a tree, unicyclic
graph, or any supplied fixed-treewidth family, the successful branch costs
$O(B)$ for fixed $w$; with $B=2/\epsppr$ this is $O(1/\epsppr)$, with no
$1/\sqrt\alpha$ factor.  The same APPR fair fallback remains unconditional.
\end{corollary}

\begin{proof}
Eliminating a bag of at most $w+1$ live variables takes $O(w^3)$ arithmetic
and creates no factor outside the certified filled graph.  There are
$O(|U|)$ bags/vertices and $|U|\leq\vol(U)\leq B$; a reverse traversal
recovers the solution with the same or smaller work.  The Green truncation
proof is unchanged.  For a tree, leaf elimination and back substitution scan
each retained incidence only a constant number of times.  A cycle has width
two.  Tree/cycle structure and the corresponding order can be verified while
the retained rows are scanned; a general fixed-width order is part of the
structural lane's certificate.  In finite precision, compute the full
principal residual after recovery and polish until its norm is at most
$\alpha\tau$; this restores the same observable terminal certificate and
adds only conditioning/precision logarithms.
\end{proof}

\paragraph{High-accuracy Green--CG experiment.}
The executable one-shot lane uses exactly
\cref{thm:two-stage-green-ball-screen}: one capped BFS, one zero-start
principal CG solve stopped by the observable residual bound, and no repeated
restricted solves.  Across six deterministic graph families, three values of
$\alpha$, and
$\epsppr\in\{10^{-4},5\cdot10^{-4},10^{-3},0.002,0.005\}$, it completes the
ball in $85$ of $90$ cases and strictly beats the better of FIFO APPR and
priority SOR in $22$ as a standalone lane.  The wins comprise $13$ stars,
six binary trees, one grid, one path, and one cycle.  After charging the
one-for-one APPR race, $15$ remain strict total-work wins: ten stars, four
binary trees, and one grid.  At $(\alpha,\epsppr)=(0.01,10^{-4})$, the
ratios are $0.103$ on the binary tree, $0.489$ on the grid, $0.882$ on the
path, and $0.945$ on the cycle.  The best overall ratio is $0.0258$ on the
star.  Every semantic error is at most $0.605\epsppr$.  The median ratio over
successful screens is nevertheless $1.95$ and the maximum is $24.85$;
therefore the fair race is algorithmically essential.  This is the first
tested set-only two-stage lane to win beyond the low-minimal-polynomial star,
but it remains an instance-adaptive rather than uniformly accelerated result.

The structural tail changes the constants and the asymptotic dependence.
On five tree/treewidth-two families with $1024$--$2047$ vertices, three values
of $\alpha$, and the same five high-accuracy targets, the exact solve spends
the full $\epsppr$ budget on Green truncation and $56$ of $75$ balls fit the
cap.  Charging the implemented leaf peeling, two-sweep cyclic-core solve,
back substitution, and one residual/materialization scan, $50$ strictly beat
the direct baseline standalone and $35$ remain wins after the fair APPR
charge; their median successful ratio is $0.302$.  At
$(\alpha,\epsppr)=(0.01,10^{-4})$, the path, cycle,
star, and binary-tree ratios are respectively $0.0469$, $0.0583$, $0.0526$,
and $0.0951$.  These are structural fixed-operator wins, not evidence that
general sparse elimination is local or fill-free.

A literal APPR-coupled variant spends the same Stage~I exposure budget on a
priority-APPR prefix and uses that checkpoint as the fixed-face CG warm start.
Charging both BFS and APPR work, it wins $19$ of the $90$ cases, has median
successful ratio $1.61$, and reduces the CG count by a median four iterations.
It retains wins on all five graph families; at
$(\alpha,\epsppr)=(0.01,10^{-4})$ its ratios are $0.217$ on the binary tree,
$0.421$ on the grid, $0.782$ on the path, and $0.870$ on the cycle.  Warm
starting can increase CG iterations when the remaining residual has an
unfavorable spectral mix, so a robust fixed-face implementation may race the
zero and warm starts.  Either output uses the same residual certificate; this
choice affects work only.

A one-ball oracle diagnostic improves the standalone win count only from
$22$ to $23$.
The additional case is the high-accuracy broom: stopping at radius $47$, just
before its high-degree hub, passes the leakage gate with ratio $0.738$;
the universal Green radius crosses the hub and retains the full broom.  This
suggests one narrow implementable refinement rather than an unrestricted
multi-radius search.  The tested rule permits one leakage attempt before the
first BFS layer whose maximum degree is at least eight times the retained
median degree (and at least eight).  A failed attempt warm-starts the final
Green ball.  Exactly one of $90$ cases triggers, it passes on the broom, and
the executable standalone lane realizes all $23$ oracle wins.  Its extra
broom ratio is $0.738$, so it does not increase the $15$ fully charged fair
wins.  The threshold eight is a
declared policy, not a theorem about optimal pruning; correctness follows
entirely from the leakage certificate and the hard-cap/fallback discipline.

\subsection{APPR: the unconditional containing-envelope lane}

The provider proposition
\path{prop:aesp-cd-point-source-appr-envelope} gives
\begin{equation}
 S^\star(\rho)
 \subseteq E_{\rm APPR}(\rho)
 \subseteq S^\star\!\left(\frac{1-\alpha}{2}\rho\right),
 \qquad
 \vol(E_{\rm APPR})\leq\frac{2}{(1-\alpha)\rho}.
 \label{eq:two-stage-appr-envelope}
\end{equation}
For $\alpha<1/2$, the volume is below $4/\rho$.  Combining
\cref{eq:two-stage-appr-envelope} with
\cref{thm:two-stage-composition} yields a completely unconditional
two-stage solver, but its worst-case discovery work remains
$\Theta(1/(\alpha\rho))$.  A center-seeded star already has radius one and
envelope volume $\Theta(1/\rho)$ while forcing that APPR work.  Hence neither
small radius nor output-sized volume pays for literal repeated pushes.

Moreover, APPR's numerical vector at threshold $\alpha\epsppr$ already meets
the semantic target by the same residual comparison.  Thus this lane is a
valuable correctness fallback and proves that a suitable envelope always
exists, but its optional Stage~II is again redundant at equal accuracy.  It
is not the desired complexity breakthrough.

\subsection{Exact positive-boundary batches: a complete baseline}

There is a particularly simple literal AESP-CD Stage~I that is correct on
every Stieltjes face, although not uniformly fast.

\begin{proposition}[Exact batch-refit discovery]
\label{prop:two-stage-exact-batch-refit}
Start from the source singleton whenever its shifted load is positive.
At round $j$, solve the current principal RPPR equations exactly on $A_j$,
compute every exterior KKT key, and set
\[
 A_{j+1}=A_j\cup\{i\notin A_j:g_i(x^{A_j})>0\}.
\]
If the source load is nonpositive, return the zero solution.  Otherwise the
faces remain interior and nested inside $S^\star$, and the first round with no
positive exterior key returns exactly $S^\star$ and $x_\rho^\star$.  If a
solve-plus-scan on $A$ costs $F(A)$, the charged discovery work is
\[
 \sum_j F(A_j).
\]
Racing this interruptible lane one-for-one with direct APPR gives a complete
algorithm within factor two of direct APPR, while retaining any earlier exact
support win.
\end{proposition}

\begin{proof}
Pad the exact face point by zero and write its residual as
$g=c_\rho-Qx^{A_j}$.  It vanishes on $A_j$ and is strictly positive on the
new batch $P_j$.  On $B=A_j\cup P_j$,
\[
 x^B-(x^{A_j},0)=Q_B^{-1}g_B\geq0,
\]
because a principal Stieltjes inverse is nonnegative.  Every new connected
coordinate is strictly positive, old coordinates do not decrease, and the
positive-key support-safety argument gives $B\subseteq S^\star$.  The process
therefore terminates in finitely many admissions.  At a terminal face the
active equations and all exterior nonpositive keys are precisely the global
obstacle KKT conditions.  The scheduling claim is
\cref{thm:two-stage-certificate-race}.
\end{proof}

With the exact-tail choice $\rho=\epsppr$, the structural implementation
certifies $75$ of the $90$ benchmark cases.  Its exact lane finishes before
direct APPR in $43$; after charging the one-for-one race, $13$ nontrivial wins
remain when the exact RPPR point returns immediately.  With a separate fresh
ordinary-PPR Stage~II, $33$ standalone and $8$ fair wins remain.  The batch word
still explains the limitation: a star activates $1023$ leaves in one round
and a low-$\alpha$ binary tree doubles its batch geometrically, whereas a
representative path admits one vertex in each of $38$ consecutive rounds.
Thus exact batch refits are a safe portfolio lane and a useful
baseline, not a universal product-work proof.
Retaining the last factors changes those literal counts to $37$ standalone
and $9$ fair wins.

\subsection{Exact active-set response on structural graph classes}

On a point-source tree, scalar Schur responses and propagated activation
thresholds name a genuinely positive frontier row without probing every
rejected row.  The provider theorem
\path{thm:aesp-cd-tree-singleton-threshold} gives
\begin{equation}
 \widetilde O\!\left(
  \frac{1+\vol(S^\star)}{\sqrt\alpha}
 \right)
 \leq
 \widetilde O\!\left(\frac1{\rho\sqrt\alpha}\right).
 \label{eq:two-stage-tree-discovery}
\end{equation}
The unicyclic theorem and the cactus reporter hierarchy extend the same
principle by maintaining low-dimensional block responses.  These are true
Stage~I breakthroughs.  Their exact restricted RPPR values are already
$\rho$-accurate for PPR and may return immediately; if an implementation
exports only the certified face, \cref{thm:two-stage-composition} adds a
product-scale fixed-face solve.

We implemented the tree reporter directly, including the handled child
threshold heaps, lazy rooting on first active-row exposure, root-path response
updates, and one final leaf solve.  It performs no preliminary full-tree scan;
each admitted adjacency list is read once and inactive descendants are known
only through degree/load metadata.  On $60$ path, star, binary-tree, and broom cases
with $1024$--$2047$ vertices, all $60$ traces certify the exact face.  Among
the $51$ cases in which direct APPR performs nonzero work, returning the
already accurate RPPR point wins $27$ standalone and $14$ after charging a
literal one-for-one APPR race.  The strict set-only version has the same
counts: the threshold messages themselves certify the face, so Stage~II's
ordinary-PPR leaf solve replaces, rather than duplicates, recovery of the
RPPR point.  Because that solve is exact, the tuned split spends the whole
semantic budget on screening, $\rho=\epsppr$, and assigns zero terminal
budget.  Its best ratio is $0.059$ on a star; on the low-$\alpha$ path and
broom it uses $1630$ charged work versus $6312$ for direct APPR.  Every
resulting ordinary-PPR point passes the semantic error target.  The benchmark
also performs a separate RPPR solve and KKT scan as an uncharged independent
audit.  An additional $1890$ randomized small-tree runs all passed the
independent structural solve and global KKT verifier.  Thus
\cref{eq:two-stage-tree-discovery} is now both a proved structural closure and
an executable two-stage algorithm, rather than an oracle assumption.

\subsection{The general dynamic-obstacle contract}

The provider theorem
\path{thm:aesp-cd-hard-cap-dynamic-oracle} maintains a retained envelope
$U$ with
\begin{equation}
 \vol(U)\leq\frac2\rho.
 \label{eq:two-stage-hard-cap}
\end{equation}
Every positive boundary key is support-safe.  Speculative inactive rows use a
pessimistic $1/\rho$ volume ledger, and a boundary with no positive key is the
global obstacle KKT certificate.  If a dynamic obstacle/reporter oracle costs
$W_{\rm DS}(2/\rho)$, Stage~I costs
\begin{equation}
 W_{\rm disc}=W_{\rm DS}(2/\rho)+O(1/\rho).
 \label{eq:two-stage-ds-work}
\end{equation}
An oracle bound
$W_{\rm DS}(B)=\widetilde O(B/\sqrt\alpha)$ would close the arbitrary-graph
target immediately through \cref{thm:two-stage-composition}.  This is the
minimal unresolved theorem in the two-stage presentation.

\subsection{AESP-CD as a Stage I primitive}

On any fixed retained envelope $U$, accelerated projected-gradient scratch
plus the safe obstacle retraction solves the restricted problem in
\begin{equation}
 \widetilde O(\vol(U)/\sqrt\alpha)
 \label{eq:two-stage-obstacle-primitive}
\end{equation}
work without a strict active or dual margin; this is the provider result
\texttt{cor:aesp-cd-margin-free-}\allowbreak
\texttt{obstacle-primitive}.  Thus the inner
restricted solve is not the missing piece.  It may either publish certified
boundary violations or invoke \cref{prop:two-stage-gap-to-envelope} once its
global gap certificate is small enough.  The missing work is obtaining that
certificate while charging exposure, boundary refresh, and response reuse
through a sequence of admissions.
Restarting a full solve after every singleton can cost quadratic cumulative
volume on a path.

The safe use of AESP-CD in Stage~I is consequently:
\begin{enumerate}[leftmargin=2.2em]
\item lock one retained envelope;
\item run signed accelerated scratch only inside it;
\item publish a lower-safe obstacle approximation;
\item certify a batch of boundary violations, or certify the global gap and
  threshold the published point into an approximate envelope;
\item discard the polynomial history before changing the envelope.
\end{enumerate}
This avoids changing-face momentum errors, but it still requires batching or
a dynamic response/reporter to achieve target total work.

\subsection{SOR and APPR inside discovery}

Gauss--Seidel, equivalently SOR with $\omega=1$, is attractive during support
discovery because zero-start obstacle updates are monotone and preserve lower
certificates.  It is a safe verifier and fallback.  Over-relaxation
$\omega>1$ may overshoot the restricted solution and therefore is not a
default support-safe publication rule; it requires an a posteriori
Stieltjes retraction or must be postponed to Stage~II.

Exact support certification from an approximate SOR state also needs a
nonzero active/dual margin or an exact fallback.  With certified margins, the
provider result \texttt{cor:aesp-cd-margin-}\allowbreak
\texttt{obstacle-primitive} identifies the
face after an accelerated fixed-envelope solve, with the margins entering
only logarithmically.  Without margins, the gap-to-envelope certificate is
the cleaner interface: SOR need only expose a lower-safe sparse state and a
valid global gap bound.
